% Pursuing Stacks (cotes 134-2 to 134-8): archive cote 134-2, site page 57 (PDF page 58).
% appendix (letters to Breen). Draft compared: ps1.tex:1387-1414.
% AI-transcribed by Claude Opus 5.5 (claude-opus-5-5, adaptive thinking, effort high), Message Batches API, batch ps1 (msgbatch_01Rv5hCaNTtUt4JHx1yqoHPL), 2026-11-12; not yet checked by eye.
\scanpage{134-2/57}
\mspage{-\over{1}{4}2 -}
\margin{(26)}
On the other hand, I am still intrigued by the following question: is
there an analogue of the ``tapis'' of Dold-Puppe\inserted{(26)} for semi-simplicial groups (\uline{not
necessarily commutative}) and what form should it take? To tell the truth I
consider the yoga
\[ \hand{(*)} \qquad \text{simplicial sets} \longleftrightarrow \infty\text{-groupoids} \]
as being essentially the ultimate ``set theoretic'' version of Dold-Puppe, which
I would deduce \inserted{from (*)} by making explicit solely the fact that the abelian groups in
($\infty$-cat) are ``nothing else'' than the \inserted{chain} complexes \deleted{of chains} in (Ab). One should
therefore \inserted{first} determine what should be the groups in ($\infty$-cat). I can tell you
what these are in (1-cat) - this will be discussed at length in the book of
Mme. Sinh, I think in the chapter ``\emph{strict} Gr-categories'' (i.e. the isomorphisms
of associativity, for unity and inverse $X.X^{-1} \simeq 1$ are \emph{identities}). One can
make explicit for example how (\emph{via} the fact that a Gr-category is Gr-equivalent
to a strict Gr-category) the calculation with the Gr-categories reduces to a
very algebraic calculation with the \emph{strict} Gr-categories, by a kind of
``calculus of fractions'' (by choice, left or right) of the type which is used
in giving the construction of derived categories. In any case, here is the
explicit formulation of the structure (groups in (1-cat)) in terms of the
\margin{\unsure{x}(27)}
theory of groups (1-categories in (Gr)).\over{\unsure{(*)}}{(27)}The structure is described by a
quadruplet $(L_1, L_0, d, \theta)$ with
\[ L_1 \xrightarrow{d} L_0 \]
a homomorphism of ordinary groups,
\[ \theta : L_0 \longrightarrow \mathrm{Aut}_{Gr}(L_1) \]
an operation of $L_0$ on $L_1$ , with the following two axioms:

(a) $d$ commutes with the operation of $L_0$ , when $L_0$ acts on $L_1$
via $\theta$ and on itself by inner automorphisms:
\[ d(\theta(x_0).x_1) = \mathrm{int}(x_0)\, d(x_1) \]
(b) \[ \theta(d(x_1) = \mathrm{int}(x_1) \ . \]
These properties imply that $\mathrm{Im}\, d$ is normal in $L_0$\over{,}{\unsure{;}} (hence $\mathrm{Coker}\, d = \pi_0$
is defined) and $\pi_1 = \mathrm{Ker}\, d$ is central in $L_1$ , and finally that $L_0$

\margin{57}
